\section{Source transformations with exact control targets}
\label{sec:source-rewrite-lemmas}

This section records ordinary universal lemmas for the Wade register-language
semantics used in \texttt{RegisterPath.lean} and \texttt{RMDag.lean}.  They are
independent of finite register testing.  The delivered 236-state candidate
uses an ordinary reassociation of three closed blocks; the jump-adjusting
rules below are a larger search space for subsequent work.

Write $(A,B)$ for a source cons cell.  A leaf is an increment $I_r$, a
decrement $D_r$, a break $B_k$, a continue $C_k$, or halt.  Decrement reduces
a positive register and continues normally; on zero it acts as a break of
level one.  All explicit break and continue operands in this section are
strictly positive.  The source frontend enforces that restriction; a checker
that interprets a zero operand by the positive-operand formulas below would
be unsound.

\subsection*{A finite control graph is a universal certificate}

Unfold a finite acyclic source DAG and number its leaves from left to right.
For every subtree prefix $p$, let $[\ell(p),r(p))$ be its leaf interval.
At a leaf with path $p$, increment and positive decrement target the next
leaf, or the normal exit if it is last.  Zero decrement uses the break-one
target.  A break of level $k$ finds the $k$th zero in $p$, counted from the
right, and targets the end of the subtree preceding that zero.  A continue
finds the $k$th one and targets the beginning of its preceding subtree.
If fewer than $k$ matching bits exist inside the local subtree, retain the
external exit type and the number of ancestors still required.

More explicitly, a local control target is a leaf ordinal, a halt outcome,
or an external token $\mathsf B_j$ or $\mathsf C_j$.  The token
$\mathsf B_0$ denotes normal completion, while positive subscripts denote
unconsumed break or continue levels.  For $(A,B)$ with $m$ leaves in $A$,
the rules that embed local targets are
\[
\begin{array}{c|cc}
\text{Local target}&\text{from }A&\text{from }B\\ \hline
i&i&m+i\\
\mathsf B_0&m&\mathsf B_0\\
\mathsf B_j\ (j\ge1)&\mathsf B_{j-1}&\mathsf B_j\\
\mathsf C_1&\mathsf C_1&0\\
\mathsf C_j\ (j\ge2)&\mathsf C_j&\mathsf C_{j-1}\\
\mathsf{halt}&\mathsf{halt}&\mathsf{halt}
\end{array}
\tag{C}
\]
Here an internal $i$ in the first or second column is relative to the
corresponding child.  In the complete well-formed program, normal completion
restarts the root at leaf zero, and no positive external token remains.

\paragraph{Certificate lemma.}
Suppose two finite source trees have the same number of leaves.  At each
corresponding ordinal, require the same register action and the same ordered
control targets, including all external tokens.  For a decrement, the zero
and positive targets must agree separately.  Then the two trees have the
same execution for every initial register file when configurations are
identified by current leaf ordinal and register file.  They can be substituted
for one another in every surrounding source context.

To prove this, note first that \texttt{RegisterPath.execute} uses the path
truncated and zero-extended to precisely the depth of the decoded leaf.
Thus the next decoded leaf depends only on that leaf and, for decrement,
whether the selected register is zero.  The preceding target formulas are
therefore exact.  Matching action and target data gives a one-step graph
isomorphism, and induction gives the result for every finite execution;
determinism extends it to divergence.  The compositional rules (C) prove
context substitution by induction on the surrounding tree.  They also
explain why residual external levels must be compared, rather than merely
noting that both jumps leave the subtree.

An unconditional break or continue has no register effect.  Its numeric
operand is representation data once its target is known.  Consequently an
equivalence checker may ignore that operand after proving exact target
agreement.  Even its break/continue opcode can be erased in a full-program
graph comparison if the resulting targets, including root wraparound, are
equal.  Preserving the opcode is a sound, more restrictive search rule.

\subsection*{Reassociation with adjusted jumps}

Consider
\[
             (A,(B,C))\quad\longmapsto\quad((A',B),C').
\tag{L}
\]
The paths of leaves in the three blocks change from
$0\alpha,10\beta,11\gamma$ to $00\alpha,01\beta,1\gamma$.
For a jump inside a block, call it \emph{escaping} when its operand exceeds
the number of matching zero or one bits in its path relative to that block.
The difference is its residual level.

The following sufficient conditions and changes preserve the exact local
control graph in (L):
\begin{enumerate}
\item No zero-decrement branch escapes $A$.
\item Increase by one every explicit break operand escaping $A$; these
changes define $A'$.  Leave its other leaves unchanged.
\item Every continue escaping $C$ has residual level at least two.  Decrease
each such operand by one; these changes define $C'$.
\item Leave all internal jumps, register operands, and all leaves of $B$
unchanged.
\end{enumerate}

For a direct proof, internal targets are unchanged by the prefix replacement.
An escaping break from $A$ crosses one left edge before rotation and two
afterward, hence the increment of its operand.  Its implicit decrement-zero
branch cannot receive that increment, which explains condition 1.  Escaping
continues from $C$ cross two right edges before rotation and one afterward,
hence the decrement.  A residual continue of level one from $C$ would restart
$B$ before rotation; after rotation no ancestor of $C$ begins at $B$.
This target is unrepresentable by changing only that continue operand, which
explains condition 3.  In $B$, the outer edge pair changes from $10$ to $01$:
a break crosses one left edge in either form, and a continue crosses one
right edge in either form.  Residual continue one restarts $A$ in both, and
residual break one ends the whole displayed block in both.  All other exits
follow immediately from (C).

The inverse rule decreases explicit break operands escaping $A$ by one,
requires their old residual levels to be at least two and forbids an escaping
implicit zero-decrement branch there, and increases each continue operand
escaping $C$ by one.  No other adjustment is needed.

For example, without any register inequality assumption,
\[
 (I_a,(I_b,C_{k+1}))\ \simeq\ ((I_a,I_b),C_k),\qquad k\ge1.
\tag{I}
\]
For $k=1$ the whole block restarts on both sides.  For $k\ge2$ it increments
the same two registers and emits the same residual continue of level
$k-1$.  Formula (I) cannot be extended by retaining $C_1$ on the left and
inventing $C_0$ on the right: the left program repeats $I_b$ without repeating
$I_a$, and zero-continue is outside the stated source language.

These rules need not be implemented as ad hoc syntactic exceptions.  Given
any new binary parenthesization with the same ordered leaves, recompute its
available target for each positive jump operand and choose an operand that
matches the old target.  Check every decrement-zero edge separately because
its level is fixed.  If all targets can be matched, the certificate lemma
proves the transformation.  The new parenthesization may then be hash-consed
to expose source DAG sharing.  Its emitted Turing-machine count must still
be measured independently.

\subsection*{When increment commutation is valid}

For a literal two-leaf subtree, the identity
\[
                    (I_a,I_b)\simeq(I_b,I_a)
\tag{S}
\]
holds in every surrounding context.  Both terminate normally with the same
register file; it also holds when $a=b$.  More generally, any permutation of
a finite, always-executed increment sequence has the same net addition
vector.

For consecutive leaves in a larger control graph, adjacency alone is
insufficient.  A sufficient graph condition is that the sequence has one
entry: its internal leaves have no incoming edge except the predecessor's
fallthrough, and the initial program entry is not an internal leaf.  The
sequence can then be treated as one block with the same output register file.
Without that condition, an outside continue may enter the second increment
and observe which register is incremented there.  Combining (I) with (S),
when its jump condition holds, safely creates a commutable pair that was
previously separated by the source tree's right nesting.

\subsection*{Moving an increment through a transfer loop}

For $r\ne s$, define the closed transfer block
\[
 T_{r,s}=(D_r,(I_s,C_2)).
\]
On every register file $R$ it terminates normally, sets register $r$ to zero,
sets $s$ to $R(s)+R(r)$, and leaves all other registers unchanged.  This is
the existing generic \texttt{RMDag.transfer} contract.  Induction on $R(r)$
also proves it directly.

We obtain the three-way identity
\[
 (I_r,T_{r,s})\ \simeq\ (I_s,T_{r,s})
 \ \simeq\ (I_s,(D_r,C_2)),\qquad r\ne s.
\tag{T}
\]
All three terminate normally with
\[
 R'(r)=0,\qquad R'(s)=R(s)+R(r)+1,\qquad
 R'(j)=R(j)\quad(j\ne r,s).
\]
The final form executes $I_s$ before each test, including the final zero test;
there are $R(r)+1$ increments.  This proves the displayed normal form without
a register invariant.  Determinism then gives equality of every exit/output
judgment, so the replacement is valid in arbitrary source contexts.

The first equality in (T) also applies to a nonlocal control-graph pattern.
Let the first increment fall through to a decrement of the same register;
require the decrement's positive edge to go to $I_s$, that increment to an
unconditional continue, and that continue back to the same decrement.  Its
zero edge may be any unchanged continuation.  Starting before the first
increment, the old and new executions both reach that continuation with the
displayed output.  Other edges may enter the loop, because none of its
instructions is changed by this first equality.

Several such increment changes may be made together when no changed
increment is part of another witness's loop body.  This disjointness of the
changed sites from the witness interiors makes every witness survive the
other changes.  The measured file
\texttt{reference/transfer\_source\_receipt.json} contains 34 exact witnesses,
and the set of their first increment sites has empty intersection with the
set of their loop increment sites.  The source theorem therefore justifies
the simultaneous changes in that particular experiment.  It does not by
itself claim any reduction in emitted Turing-machine states.

\paragraph{Proof status.}
The finite-control certificate and transformations above are ordinary
mathematical consequences of the stated semantics.  No new compiled Lean
theorem or internal PA proof is claimed here.  A concrete source candidate
also needs its finite unfolding/control check and the complete source-to-TM
calibration; a smaller source DAG alone is not a state bound.
